\chapter{Veredicto Pericial y Síntesis del Caso}

La fase final de toda investigación pericial consiste en integrar la totalidad de los indicios materiales, marcas de tiempo y correlaciones de red en una hipótesis concluyente que reconstruya la mecánica de los hechos y determine el grado de afectación real de la infraestructura evaluada.

\section{Doble Vector de Intrusión y Perfil de los Atacantes}

El análisis forense integral demuestra con certeza pericial la convergencia de \textbf{dos intrusiones de naturaleza e intencionalidad claramente diferenciadas}, comparadas en el Cuadro~\ref{tab:comparativa_intrusiones}:

\begin{table}[htbp]
\centering
\small
\begin{tabular}{lp{6.5cm}p{6.5cm}}
\toprule
\textbf{Dimensión Pericial} & \textbf{Intrusión 1: Oportunista Externa} & \textbf{Intrusión 2: Espionaje Interno} \\
\midrule
\textbf{Vector de Entrada} & Fuerza bruta FTP / Cuentas con claves débiles (\texttt{cgarcia}, \texttt{alopez}) & Acceso a consola local / Aprovechamiento del desconcierto \\
\textbf{Móvil Principal} & Apropiación de ancho de banda y CPU para bots IRC (\textit{psyBNC}, \textit{Eggdrop}) & Espionaje de contraseñas de \texttt{root}, intercepción de correo y control nodal \\
\textbf{Herramientas Clave} & \texttt{in.ftpd}, \texttt{emech}, \texttt{psybnc}, scripts de retransmisión & Rootkit Suckit (\texttt{sk}, \texttt{init!}), sniffer \texttt{in.telnetd}, \texttt{clean!} \\
\textbf{Evasión Empleada} & Directorios ocultos (\texttt{...}, \texttt{.king}) & Parcheo de kernel vía \texttt{/dev/kmem}, atributos \texttt{chattr +i}, sufijo \texttt{!} \\
\textbf{Perfil Técnico} & Automatizado, oportunista, dispersión geográfica internacional & Usuario con acceso físico al edificio, conocimiento de topología, torpeza operativa \\
\bottomrule
\end{tabular}
\caption{Matriz comparativa de los dos vectores de intrusión convergentes.}
\label{tab:comparativa_intrusiones}
\end{table}

\begin{enumerate}
    \item \textbf{La intrusión oportunista externa (servicios parásitos):}
    El punto de acceso inicial a la infraestructura fue el servicio FTP, complementado con intentos de explotación sobre Samba y montajes NFS. La causa determinante de este acceso fue la debilidad en la política de contraseñas de los usuarios locales; se identificó que las cuentas \texttt{cgarcia}, \texttt{alopez} y \texttt{jvaldez} fueron vulneradas mediante ataques de fuerza bruta basados en diccionarios debido a la simplicidad de sus claves.
    
    El objetivo principal de este actor externo fue el despliegue de servidores IRC clandestinos y bots de retransmisión para apropiarse del ancho de banda y la capacidad de cómputo del servidor. Este tipo de atacantes busca preservar los servicios activos pasando inadvertidos, sin sustraer bases de datos ni causar daño destructivo a los sistemas operativos.

    \item \textbf{La intrusión interna y espionaje administrativo:}
    Los artefactos localizados en el directorio \texttt{/dev} revelan un segundo ataque, cronológicamente independiente y posterior al despliegue de los servicios IRC. El móvil fundamental de este segundo evento fue la obtención del control absoluto del equipo, la interceptación de las contraseñas de \texttt{root} en caso de que fuesen reemplazadas por el personal de cómputo, y la captura de información confidencial de las cuentas de usuario mediante \textit{sniffers} locales.
\end{enumerate}

La coexistencia temporal y la divergencia metodológica entre ambos ataques se esquematizan en la Figura~\ref{fig:doble_vector_ataque}.

\begin{figure}[htbp]
\centering
\begin{tikzpicture}[
    scale=0.88,
    every node/.style={transform shape, font=\sffamily},
    extstep/.style={draw=secondary, thick, fill=secondary!10, rounded corners=3pt, align=center, font=\scriptsize, minimum width=2.5cm, minimum height=1.1cm},
    intstep/.style={draw=accent, thick, fill=accent!10, rounded corners=3pt, align=center, font=\scriptsize, minimum width=2.5cm, minimum height=1.1cm},
    flow/.style={draw=secondary, ultra thick, -{Stealth[length=2.5mm]}},
    intflow/.style={draw=accent, ultra thick, -{Stealth[length=2.5mm]}}
]

% Vector 1: Oportunista Externo
\node[extstep] (e1) at (0, 3.2) {\textbf{Fuerza Bruta FTP}\\\tiny Cuentas débiles (\texttt{proa})};
\node[extstep] (e2) at (3.5, 3.2) {\textbf{Sembrado IRC}\\\tiny \texttt{psybnc}, \texttt{emech}};
\node[extstep] (e3) at (7.0, 3.2) {\textbf{Canal C2}\\\tiny Conexión UnderNET 6667};
\node[extstep] (e4) at (10.5, 3.2) {\textbf{Objetivo Logrado}\\\tiny Parásito de CPU / Tráfico};

\draw[flow] (e1) -- (e2);
\draw[flow] (e2) -- (e3);
\draw[flow] (e3) -- (e4);

% Vector 2: Espionaje Interno
\node[intstep] (i1) at (0, 0.4) {\textbf{Acceso Consola Local}\\\tiny Aprovecha confusión inicial};
\node[intstep] (i2) at (3.5, 0.4) {\textbf{Suckit Rootkit}\\\tiny Parcheo \texttt{/dev/kmem} y \texttt{init!}};
\node[intstep] (i3) at (7.0, 0.4) {\textbf{Sniffer de Terminal}\\\tiny Intercepción en \texttt{in.telnetd}};
\node[intstep] (i4) at (10.5, 0.4) {\textbf{Objetivo Logrado}\\\tiny Captura de clave \texttt{root} y correo};

\draw[intflow] (i1) -- (i2);
\draw[intflow] (i2) -- (i3);
\draw[intflow] (i3) -- (i4);

% Etiquetas de vectores
\node[font=\bfseries\small\color{secondary}, anchor=south west] at (-0.2, 3.9) {Vector 1: Intrusión Oportunista Externa (Bots IRC / Sin interés en datos locales)};
\node[font=\bfseries\small\color{accent}, anchor=south west] at (-0.2, 1.1) {Vector 2: Intrusión y Espionaje Interno (Captura de credenciales administrativas / Lectura de correo)};

% Línea divisoria / correlación
\draw[draw=gray!40, dashed, thick] (-0.3, 2.0) -- (11.8, 2.0) node[pos=0.5, above, font=\tiny\itshape\color{gray}] {Intervención concurrente e independiente en el mismo servidor (\texttt{helios})};

\end{tikzpicture}
\caption{Esquema pericial de los dos ciclos de ataque independientes identificados en la investigación.}
\label{fig:doble_vector_ataque}
\end{figure}

El inicio de este segundo vector coincide temporalmente con la sustitución preventiva de contraseñas de administración en todas las máquinas departamentales. Fue inmediatamente después de dicho cambio cuando aparecieron las primeras evidencias de espionaje sobre el tráfico de red. La evidencia sugiere que el hallazgo inicial de los servicios IRC alentó a un usuario local a ejecutar sus propias acciones maliciosas, amparándose en la confusión generada por la intrusión externa preliminar.

La impericia técnica en la instalación del \textit{rootkit} ---reflejada en errores al fijar atributos de inmutabilidad y la omisión de modificar los propietarios originales de los archivos bajo \texttt{/dev}--- indica un atacante con conocimientos teóricos básicos de \textit{cracking}, muy distante de la sofisticación de los operadores externos tradicionales.

\section{Evaluación de Responsabilidad y Atribución Local}

La investigación de los servidores bajo administración departamental demostró que el interés del intruso se concentró exclusivamente en \texttt{helios}, \texttt{atlas} e \texttt{io}, las máquinas nodales del centro de cómputo. El atacante conocía con precisión la arquitectura de hardware, la función operativa y la ubicación física exacta de los gabinetes dentro del campus:
\begin{itemize}
    \item Las auditorías sobre los conmutadores y subredes adyacentes demostraron que no existieron ataques ni accesos anómalos contra sistemas ubicados fuera del edificio departamental.
    \item Un actor remoto externo carece de los medios para inferir la distribución física y la relación de confianza funcional que unía a estos servidores específicos.
\end{itemize}

El patrón de actividad registrado en la cuenta \texttt{duende} evidencia un conocimiento directo del compromiso del sistema y la tenencia de múltiples herramientas ofensivas empleadas en el ataque. La ejecución de estos binarios, la inspección no autorizada de correspondencia electrónica en el servidor de correo y el intento deliberado de ocultar la evidencia modificando la propiedad de los archivos de \texttt{duende} a \texttt{root} comprometen formalmente la responsabilidad de dicha cuenta.

\begin{notaforense}[El Estándar Jurídico de la Atribución: Cuenta vs. Persona Física]
En el derecho procesal y la práctica forense judicial, existe una distinción capital entre demostrar la autoría de una \textbf{cuenta digital} y atribuir el hecho a una \textbf{persona física}:
\begin{itemize}
    \item La evidencia pericial sobre \texttt{helios}, \texttt{atlas} e \texttt{io} demostró de manera directa e irrefutable que las acciones maliciosas fueron ejecutadas por la cuenta \texttt{duende} y desde direcciones IP asignadas a conexiones locales de Morelia.
    \item Para trasladar la responsabilidad pericial hacia la persona titular de la cuenta en un proceso judicial, el perito debe respaldarse en la \textit{prueba indiciaria concurrente}: correlación de horarios de presencia física en el laboratorio, titularidad de la línea telefónica del ISP certificada por orden judicial, testimonios y ausencia de indicios de usurpación o malware en el equipo personal del sospechoso.
\end{itemize}
\end{notaforense}

\section{Medidas de Contención, Remediación y Cierre Pericial}

Como parte de la estrategia de contención y saneamiento:
\begin{itemize}
    \item Los servidores afectados fueron retirados de producción y sustituidos por equipos nuevos desde el inicio de los exámenes \textit{post-mortem}. Tras dicho reemplazo, cesaron de manera definitiva los intentos de conexión anómalos contra la infraestructura, un patrón incompatible con escaneos externos desasociados.
    \item Se implementaron señuelos periciales (\textit{honeypots}) para monitorizar potenciales reincidencias sin comprometer datos reales.
    \item Las instituciones y redes externas cuyos equipos figuraban comprometidos como nodos de retransmisión fueron notificadas a través de sus respectivos equipos de respuesta a incidentes (CERT).
\end{itemize}

Con este dictamen, la Administración del Centro de Cómputo del Instituto de Física y Matemáticas de la UMSNH dio por concluidas las indagatorias formales del incidente, tras acumular más de 180 horas de trabajo pericial de campo y análisis documental exhaustivo.

\begin{retoforense}[Defensa del Dictamen Pericial en Juicio Oral]
Si la defensa legal del sospechoso alegara en sede judicial que un intruso externo desconocido vulneró la cuenta \texttt{duende} mediante un troyano para incriminarlo falsamente (\textit{the Trojan Horse defense}), ¿qué elementos técnicos recabados en este peritaje (correlación de accesos locales, archivos de configuración de Midnight Commander, atributos de inmutabilidad erróneos y accesos simultáneos) sustentan de forma sólida la hipótesis pericial del informe?
\end{retoforense}
